\section{The Change in Game: From Getting Answers to Giving Reasons}
\chapterline

In the opening pages of this book (\emph{How This Book Should Be Read}), we distinguished between calculating numerical answers and providing deductive proofs. If your background is in engineering or applied science, your training naturally directs your attention toward methods of solution. When you encounter a system of linear equations, you reach for Gaussian elimination or matrix factorizations. When a differential equation appears, you check whether it is separable, linear, or amenable to an integrating factor. When an integral looks intimidating, you scan your mental catalog for substitutions, integration by parts, or numerical quadrature.

That computational agility is valuable. But pure mathematics asks you to pause before reaching for an algorithm. It shifts the central question from \emph{what is the answer?} to something fundamentally different:

\begin{quote}
\itshape What exact claim is being made here, and why is it logically impossible for it to fail under our assumptions?
\end{quote}

Answering that question requires a different kind of argument than computing a result.

\subsection{Checking numbers versus establishing truth}

Consider a simple numerical experiment. You add consecutive odd numbers and watch the running totals emerge:
\begin{align*}
1 &= 1 = 1^2, \\
1 + 3 &= 4 = 2^2, \\
1 + 3 + 5 &= 9 = 3^2, \\
1 + 3 + 5 + 7 &= 16 = 4^2, \\
1 + 3 + 5 + 7 + 9 &= 25 = 5^2.
\end{align*}

A striking regularity appears immediately. It is natural to form a conjecture---a proposed statement believed to be true based on patterns, but not yet proved---that for any positive integer $n$ (where $\N = \{1, 2, 3, \ldots\}$ denotes the natural numbers):
\[
1 + 3 + 5 + \cdots + (2n - 1) = n^2.
\]

You could write a quick script to test this equation for $n = 10$, $n = 1{,}000$, or even the first ten billion natural numbers. Every single test will return a match. 

In engineering, ten billion consecutive successful trials is strong empirical evidence. In pure mathematics, however, those ten billion checks leave infinitely many natural numbers untested. Checking specific numbers, no matter how many, cannot prove a claim for every natural number $n$.

This exact dilemma confronted Thales, Pythagoras, and Euclid in the sixth through third centuries BCE. Working with sand trays, straightedges, and counting pebbles, the early Greeks realized that empirical testing is inherently limited on infinite sets.

Several famous conjectures held for thousands or millions of test cases before failing at a single counterexample---a concrete instance where a general claim fails:
\begin{itemize}
\item \textbf{Fermat's Prime Conjecture (1640)}: Pierre de Fermat noticed that the numbers $F_n = 2^{2^n} + 1$ are prime (whole numbers greater than $1$ with no positive divisors other than $1$ and themselves) for $n = 0, 1, 2, 3, 4$:
\[
F_0 = 3, \quad F_1 = 5, \quad F_2 = 17, \quad F_3 = 257, \quad F_4 = 65{,}537.
\]
All five are prime. Fermat conjectured that $F_n$ is prime for \emph{every} natural number $n$. Almost a century later, in 1732, Leonhard Euler factored the very next case:
\[
F_5 = 2^{32} + 1 = 4{,}294{,}967{,}297 = 641 \times 6{,}700{,}417.
\]
In fact, no other Fermat prime has ever been discovered.
\item \textbf{Euler's Sum of Powers (1769)}: Euler conjectured that at least $k$ positive $k$-th powers are required to sum to another $k$-th power. For fourth powers, this asserted that $a^4 + b^4 + c^4 = d^4$ has no positive integer solutions. For more than two centuries, mathematicians found no counterexample, and twentieth-century computer searches turned up nothing. Yet in 1986, using elliptic curves and algebraic geometry, Noam Elkies showed that a counterexample exists:
\[
2{,}682{,}440^4 + 15{,}365{,}639^4 + 18{,}796{,}760^4 = 20{,}615{,}673^4.
\]
\item \textbf{The P\'olya Conjecture (1919)}: George P\'olya conjectured that for any integer $n > 1$, at least half of the integers up to $n$ have an odd number of prime factors (counted with repetition, so $12 = 2 \times 2 \times 3$ has three prime factors, which is odd, while $4 = 2 \times 2$ has two, which is even). The conjecture was verified for millions of integers. It holds without exception for every single number up to $906{,}150{,}256$. The smallest counterexample only appears at $n = 906{,}150{,}257$!
\end{itemize}

If mathematicians had stopped checking at nine hundred million, they would have accepted a false claim. Checking individual numbers only verifies finitely many cases; a proof shows why a general statement holds for every $n$.

\begin{recoverybox}
\textbf{Pause and think:} Suppose someone writes a script that tests an arithmetic identity for the first trillion positive integers, and every single case succeeds. In software testing or material science, this is overwhelming evidence. Why does a pure mathematician still refuse to declare the identity true?
\end{recoverybox}

\begin{center}
\begin{tikzpicture}[
    font=\small,
    ruler/.style={draw=gray!50, thick},
    panel/.style={draw=gray!30, fill=gray!2, rounded corners=4pt, inner sep=8pt}
]

% Panel A: Sampling
\node[panel, text width=13.0cm, minimum height=2.4cm] (pA) at (6.5, 3.0) {};
\node[anchor=west, font=\footnotesize\scshape, text=quotientorange!90!black] at (0.3, 4.0) {Empirical Verification: Pointwise Sampling};

% Number line for Panel A
\draw[ruler] (0.8, 2.7) -- (6.0, 2.7);
\node[font=\small, text=gray!70] at (6.4, 2.7) {$\cdots$};
\draw[ruler, -{Latex[length=3mm]}] (11.7, 2.7) -- (12.2, 2.7) node[right=4pt, font=\scriptsize] {$n \in \N$};
\foreach \x/\label in {1.4/1, 2.4/2, 3.4/3, 4.4/4, 5.4/5} {
    \draw[thick, gray!70] (\x, 2.85) -- (\x, 2.55) node[below=2pt, font=\scriptsize] {$\label$};
    \node[circle, fill=targetgreen, inner sep=1.8pt] at (\x, 2.7) {};
    \node[above=3pt, font=\tiny\bfseries, text=targetgreen!80!black] at (\x, 2.7) {\checkmark};
}

% Region of untested numbers
\draw[fill=trapred!12, draw=trapred!50, rounded corners=2pt] (7.0, 2.05) rectangle (11.6, 3.35);
\node[text=trapred!90!black, font=\scriptsize\bfseries, align=center] at (9.3, 2.7) {Uninspected Infinity ($n > 10^9$)\\[2pt]\scriptsize A single failure here overturns the claim};

% Panel B: Deductive Proof
\node[panel, text width=13.0cm, minimum height=2.4cm] (pB) at (6.5, 0.2) {};
\node[anchor=west, font=\footnotesize\scshape, text=sourceblue!90!black] at (0.3, 1.2) {Deductive Proof: Universal Enclosure};

% Enclosure boundary
\draw[draw=sourceblue!80, fill=sourceblue!10, rounded corners=4pt, thick] (1.0, -0.65) rectangle (11.8, 0.65);
\node[text=sourceblue!90!black, font=\scriptsize\bfseries, align=center] at (6.4, 0.42) {Structural Guarantee: covers every $n \in \N$ simultaneously};

\draw[ruler, -{Latex[length=3mm]}] (0.8, -0.05) -- (12.2, -0.05) node[right=4pt, font=\scriptsize] {$n \in \N$};
\foreach \x/\label in {1.4/1, 2.4/2, 3.4/3, 4.4/4, 5.4/5} {
    \draw[thick, gray!70] (\x, 0.10) -- (\x, -0.20) node[below=2pt, font=\scriptsize] {$\label$};
    \node[circle, fill=sourceblue, inner sep=1.8pt] at (\x, -0.05) {};
}
\node[font=\small, text=sourceblue!80!black] at (6.3, -0.05) {$\cdots$};
\foreach \x/\label in {7.4/k, 8.6/k+1, 9.8/k+2} {
    \draw[thick, gray!70] (\x, 0.10) -- (\x, -0.20) node[below=2pt, font=\scriptsize] {$\label$};
    \node[circle, fill=sourceblue, inner sep=1.8pt] at (\x, -0.05) {};
}
\node[font=\small, text=sourceblue!80!black] at (10.7, -0.05) {$\cdots$};

\end{tikzpicture}\\[0.4em]
{\small\textbf{Figure 1.1:} \textit{Pointwise sampling versus universal enclosure.} Empirical verification (top) tests individual numbers, leaving infinitely many uninspected values where a counterexample can hide. Deductive proof (bottom) identifies an underlying structural relationship that establishes all infinitely many cases at once.}
\end{center}

To establish that our odd-number pattern holds across the entire infinite tail in Figure~1.1, we need a single structural argument that covers every $n$ at once.

\subsection{Seeing the reason geometrically}

Why \emph{must} the sum of the first $n$ odd numbers equal $n^2$? We do not need a computer to check cases. We can arrange the numbers visually into a square.

Start with a single $1 \times 1$ block. To expand this block into a $2 \times 2$ square, you must wrap around its top and right edges. How many new unit squares are required? Exactly $3$. 

To expand the $2 \times 2$ square into a $3 \times 3$ square, you again wrap an L-shaped band around the outer border. That band consists of two sides of length $2$, plus the single corner block where they meet: $2 + 2 + 1 = 5$ unit squares.

\begin{center}
\begin{tikzpicture}[scale=0.95, font=\small]

% 1x1 block
\draw[fill=sourceblue!25, draw=sourceblue!80, thick] (0,0) rectangle (1,1);
\node at (0.5, 0.5) {$1$};

% 2x2 layer (3 blocks)
\draw[fill=targetgreen!25, draw=targetgreen!80, thick] (1,0) rectangle (2,1);
\draw[fill=targetgreen!25, draw=targetgreen!80, thick] (1,1) rectangle (2,2);
\draw[fill=targetgreen!25, draw=targetgreen!80, thick] (0,1) rectangle (1,2);
\node at (1.5, 0.5) {$+3$};
\node at (1.5, 1.5) {};
\node at (0.5, 1.5) {};

% 3x3 layer (5 blocks)
\foreach \x/\y in {2/0, 2/1, 2/2, 1/2, 0/2} {
    \draw[fill=quotientorange!25, draw=quotientorange!80, thick] (\x,\y) rectangle (\x+1,\y+1);
}
\node at (2.5, 1.5) {$+5$};

% 4x4 layer (7 blocks)
\foreach \x/\y in {3/0, 3/1, 3/2, 3/3, 2/3, 1/3, 0/3} {
    \draw[fill=structurepurple!25, draw=structurepurple!80, thick] (\x,\y) rectangle (\x+1,\y+1);
}
\node at (3.5, 1.5) {$+7$};

% Dimension annotations
\draw[|-|, thick, gray!80] (0, -0.4) -- (4, -0.4) node[midway, below=2pt] {width: $n = 4$ units};
\draw[|-|, thick, gray!80] (-0.4, 0) -- (-0.4, 4) node[midway, left=2pt] {height: $n = 4$ units};

% Legend / Explanation on the right
\node[anchor=west, text width=5.5cm, inner sep=6pt, fill=gray!3, draw=gray!25, rounded corners=3pt] at (5.2, 2.0) {
    \textbf{Geometric Gnomon Construction:}\\[0.3em]
    \footnotesize
    An $(n-1) \times (n-1)$ square already contains $(n-1)^2$ tiles.\\
    Wrapping the outer border requires:
    \[
    (n-1) + (n-1) + 1 = 2n - 1
    \]
    tiles. Each successive odd number is precisely the geometric layer needed to build the next perfect square.
};

\end{tikzpicture}\\[0.4em]
{\small\textbf{Figure 1.2:} \textit{The geometric gnomon construction.} Successive odd numbers wrap around an existing square like nested shells. Expanding an $(n-1) \times (n-1)$ square to an $n \times n$ square requires two strips of length $n-1$ plus a single corner tile: $(n-1) + (n-1) + 1 = 2n-1$ tiles. The algebraic expansion and the geometric tiling match exactly.}
\end{center}

In general, if you already have an $(n-1) \times (n-1)$ square containing $(n-1)^2$ tiles, completing the $n \times n$ square requires adding:
\[
n^2 - (n-1)^2 = n^2 - (n^2 - 2n + 1) = 2n - 1
\]
new tiles. Geometrically, this breaks into two arms of length $n-1$ along the top and right, plus $1$ single corner block where they meet:
\[
(n - 1) + (n - 1) + 1 = 2n - 1.
\]
This L-shaped border (known since ancient Greek geometry as a \emph{gnomon}, from the Greek word for a carpenter's square or indicator) shows that the $n$-th odd number is structurally tailored to convert an $(n-1)^2$ square into an $n^2$ square. 

We did not need to test individual cases up to $n = 10^{12}$. The geometric argument reveals the algebraic reason that forces the equation to hold across all natural numbers at once.

\begin{recoverybox}
\textbf{Pause and think:} Look at the gnomon equation $n^2 - (n-1)^2 = 2n - 1$. What happens if you add consecutive \emph{even} numbers instead: $2 + 4 + 6 + \cdots + 2n$? Can you picture what shape is formed if you attach one extra column of unit tiles to the side of the square? (Can you see why the sum is $n(n+1)$?)
\end{recoverybox}

\subsection{From formulas to structural invariants}

In engineering, formulas often serve as compact input-output converters:
\[
\vec{F} = m\vec{a}, \qquad e^{i\theta} = \cos\theta + i\sin\theta, \qquad A\vec{x} = \vec{b}.
\]
In pure mathematics, formulas are often consequences of broader geometric and algebraic structures. Instead of asking only what numerical output a formula produces, pure mathematics investigates the structural rules governing the systems where the formula lives:
\begin{itemize}
\item What is a \emph{set} (a well-defined collection of distinct objects), and how do we build collections without falling into self-referential paradoxes?
\item What is a \emph{function} (not merely an algebraic formula like $f(x) = x^2$, but an unambiguous assignment from each element of an input domain to an output target)?
\item What does it mean for two mathematical systems to share the ``same structure''---an \emph{isomorphism} (from the Greek \emph{isos} ``equal'' and \emph{morphe} ``form'': a structural dictionary that matches every element and operation of one system to another without loss of information)?
\item What properties remain \emph{invariant} (unchanged under transformations, such as stretching, twisting, or coordinate shifts)?
\end{itemize}

This shift toward abstraction unifies seemingly unrelated problems. A single theorem about \emph{groups}---mathematical systems capturing symmetry, defined as collections of actions that can be combined and undone under strict rules---explains properties of modular arithmetic, geometric symmetries of crystals, roots of polynomial equations, and conservation laws in physics.

\begin{historybox}
Early Greek mathematicians studied number theory by arranging pebbles (\emph{psephoi}, from which the word \emph{calculus} originates) into geometric arrays: triangles, squares, and rectangles. The L-shaped gnomon was not an abstract algebraic trick; it was a physical wooden carpenter's square placed around an existing pebble array to enlarge it.

The limitation of pebble arithmetic and empirical testing is that physical arrangements only verify finitely many instances at a time. The shift to deductive proof---exemplified by Euclid's \emph{Elements}---aimed to capture all infinitely many cases at once by analyzing the structure rather than counting stones.
\end{historybox}

\subsection{General proofs and the limits of visual intuition}

A structural argument like our gnomon tiling covers every positive integer $n$ at once, rather than checking values one by one.

Even so, visual and geometric reasoning has a real limitation: \textbf{our eyes can easily be deceived}. Drawings rely on physical ink and human visual perception, both of which have finite resolution. A slight tilt or a tiny gap can easily pass unnoticed, making false arguments appear plausible.

To see how unassisted visual intuition can be misleading, consider the classic dissection paradox shown in Figure~1.3. An $8 \times 8$ square of area $64$ is sliced into four pieces: two right triangles (each with legs $3$ and $8$, area $12$) and two trapezoids (each with parallel sides $3$ and $5$ and width $5$, area $20$). When rearranged, the pieces appear to fit together into a $13 \times 5$ rectangle of area $65$.

\begin{center}
\begin{tikzpicture}[scale=0.48, font=\small]

% Panel 1: The 8x8 Square (Area 64)
\begin{scope}[xshift=0cm]
    \node[anchor=south, font=\footnotesize\bfseries\scshape, text=gray!80!black] at (4, 8.3) {Square: $8 \times 8$, Area $= 64$};
    
    % Grid
    \draw[step=1, gray!18, very thin] (0,0) grid (8,8);
    
    % Trapezoid A (lower-left)
    \filldraw[fill=quotientorange!25, draw=quotientorange!80, thick] (0,0) -- (5,0) -- (5,3) -- (0,5) -- cycle;
    \node[text=quotientorange!90!black, font=\scriptsize\bfseries] at (2.2, 2.2) {$A$ (area 20)};
    
    % Trapezoid B (upper-left)
    \filldraw[fill=structurepurple!25, draw=structurepurple!80, thick] (0,5) -- (5,3) -- (5,8) -- (0,8) -- cycle;
    \node[text=structurepurple!90!black, font=\scriptsize\bfseries] at (2.2, 5.8) {$B$ (area 20)};
    
    % Triangle C (lower-right)
    \filldraw[fill=sourceblue!25, draw=sourceblue!80, thick] (5,0) -- (8,0) -- (8,8) -- cycle;
    \node[text=sourceblue!90!black, font=\scriptsize\bfseries] at (7.0, 2.5) {$C$ (12)};
    
    % Triangle D (upper-right)
    \filldraw[fill=targetgreen!25, draw=targetgreen!80, thick] (5,0) -- (5,8) -- (8,8) -- cycle;
    \node[text=targetgreen!90!black, font=\scriptsize\bfseries] at (6.0, 5.5) {$D$ (12)};
    
    % Outer border
    \draw[thick, gray!70] (0,0) rectangle (8,8);
    \draw[|-|, thick, gray!60] (0, -0.6) -- (8, -0.6) node[midway, below=1pt, font=\scriptsize] {width: $8$};
    \draw[|-|, thick, gray!60] (-0.6, 0) -- (-0.6, 8) node[midway, left=1pt, font=\scriptsize] {height: $8$};
\end{scope}

% Panel 2: The 13x5 Rectangle (Apparent Area 65)
\begin{scope}[xshift=10.2cm, yshift=0cm]
    \node[anchor=south, font=\footnotesize\bfseries\scshape, text=gray!80!black] at (6.5, 5.3) {Rectangle: $13 \times 5$, Apparent Area $= 65$};
    
    % Grid
    \draw[step=1, gray!18, very thin] (0,0) grid (13,5);
    
    % Triangle C (bottom-left)
    \filldraw[fill=sourceblue!25, draw=sourceblue!80, thick] (0,0) -- (8,0) -- (8,3) -- cycle;
    \node[text=sourceblue!90!black, font=\scriptsize\bfseries] at (4.5, 0.6) {$C$ (12)};
    
    % Trapezoid A (bottom-right)
    \filldraw[fill=quotientorange!25, draw=quotientorange!80, thick] (8,0) -- (13,0) -- (13,5) -- (8,3) -- cycle;
    \node[text=quotientorange!90!black, font=\scriptsize\bfseries] at (10.5, 2.2) {$A$ (area 20)};
    
    % Trapezoid B (top-left)
    \filldraw[fill=structurepurple!25, draw=structurepurple!80, thick] (0,0) -- (5,2) -- (5,5) -- (0,5) -- cycle;
    \node[text=structurepurple!90!black, font=\scriptsize\bfseries] at (2.5, 3.2) {$B$ (area 20)};
    
    % Triangle D (top-right)
    \filldraw[fill=targetgreen!25, draw=targetgreen!80, thick] (5,2) -- (13,5) -- (5,5) -- cycle;
    \node[text=targetgreen!90!black, font=\scriptsize\bfseries] at (7.0, 4.2) {$D$ (12)};
    
    % Parallelogram gap (vertices (0,0), (5,2), (13,5), (8,3))
    \filldraw[fill=trapred!65, draw=trapred!90, thick] (0,0) -- (5,2) -- (13,5) -- (8,3) -- cycle;
    
    % Outer border
    \draw[thick, gray!70] (0,0) rectangle (13,5);
    \draw[|-|, thick, gray!60] (0, -0.6) -- (13, -0.6) node[midway, below=1pt, font=\scriptsize] {width: $13$};
    \draw[|-|, thick, gray!60] (13.6, 0) -- (13.6, 5) node[midway, right=1pt, font=\scriptsize] {height: $5$};
    
    % Note on gap
    \node[anchor=north, font=\scriptsize, text=trapred!90!black] at (6.5, -1.6) {\textbf{Hidden gap:} area $1$ (slopes $\frac{3}{8} = 0.375$ vs $\frac{2}{5} = 0.400$)};
\end{scope}

\end{tikzpicture}\\[0.4em]
{\small\textbf{Figure 1.3:} \textit{The chessboard dissection paradox ($64 \ne 65$).} Slicing an $8 \times 8$ square (left) yields four pieces totaling area $64$. Reassembling them into a $13 \times 5$ rectangle (right) seems to produce area $65$. The illusion arises because the human eye cannot see the slight bend along the diagonal: the triangle's slope is $\frac{3}{8} = 0.375$, while the trapezoid's slope is $\frac{2}{5} = 0.400$. The pieces do not meet; they enclose a thin parallelogram gap of area $1$.}
\end{center}

Where did the extra square unit come from?

The human eye sees a clean diagonal line running from $(0,0)$ to $(13,5)$. But calculate the slopes of the abutting pieces:
\[
\text{slope of triangle hypotenuse} = \frac{3}{8} = 0.375, \qquad \text{slope of trapezoid edge} = \frac{2}{5} = 0.400.
\]
Because $\frac{3}{8} \ne \frac{2}{5}$, the pieces do not meet along a straight line. They bend slightly outward, leaving a long, paper-thin parallelogram gap running along the seam. The area of that hidden gap is precisely $1$ unit square:
\[
|5 \times 3 - 8 \times 2| = |15 - 16| = 1.
\]
Because the slope difference is only $0.025$, our visual system smooths over the gap and perceives a solid rectangle. If we relied on pictures alone as proofs, we would be forced to conclude that $64 = 65$.

Geometric drawings also cannot directly represent higher dimensions. We cannot sketch a ten-dimensional vector space or an infinite-dimensional space of continuous functions on a chalkboard.

When visual intuition reaches its practical boundary, mathematics turns to precise symbolic language. In Chapter~2 (\emph{Mathematical Sentences: Statements, Predicates, and Claims}), we begin building that grammar from the ground up, learning how to distinguish expressions from assertions, open conditions from true-or-false statements, and valid deductive arguments from empirical patterns.

\subsection{Exercises: cultivating three perspectives}

\begin{enumerate}[label=\textbf{\arabic*.}]
\item \textbf{Mental Check: Patterns versus Invariants.}
\begin{enumerate}[label=(\alph*)]
\item Consider the polynomial $P(n) = n^2 - n + 41$. Compute $P(1), P(2), P(3), P(4)$. Are all four values prime? (In fact, $P(n)$ produces prime numbers for every integer from $n = 1$ up to $n = 40$.) What happens at $n = 41$? Why does this reinforce the lesson of Figure~1.1?
\item Someone claims that $(n-1)(n-2)(n-3) = 0$ is a universal law for all natural numbers, citing $n = 1, 2, 3$ as evidence. At what value of $n$ does this claim break, and what went wrong with the tester's reasoning?
\end{enumerate}

\item \textbf{Core Practice: Three Perspectives on Odd Numbers.}
For the identity $1 + 3 + 5 + \cdots + (2n-1) = n^2$, write down three distinct explanations:
\begin{enumerate}[label=(\alph*)]
\item A concrete numerical verification for $n = 6$.
\item A paragraph in plain English explaining the geometric gnomon diagram (Figure~1.2) to someone who has never studied algebra.
\item An inductive step: assuming the identity holds for some positive integer $k$, so that $1 + 3 + \cdots + (2k-1) = k^2$, show algebraically how adding the next odd number, $2(k+1) - 1$, produces $(k+1)^2$.
\end{enumerate}

\item \textbf{Core Practice: Summing Consecutive Even Numbers.}
Consider the sum of the first $n$ consecutive even numbers: $2 + 4 + 6 + \cdots + 2n$.
\begin{enumerate}[label=(\alph*)]
\item Calculate the sum for $n = 1, 2, 3, 4$. What pattern emerges in relation to $n(n+1)$?
\item By pairing each odd number in $1 + 3 + 5 + \cdots + (2n-1)$ with $1$ extra unit, deduce geometrically why the even sum forms an $n \times (n+1)$ rectangle of area $n(n+1)$.
\end{enumerate}

\item \textbf{Deeper Reflection: The Anatomy of a Visual Illusion.}
In the chessboard dissection paradox (Figure~1.3), the dimensions $3, 5, 8, 13$ are consecutive terms of the Fibonacci sequence ($F_4 = 3, F_5 = 5, F_6 = 8, F_7 = 13$).
\begin{enumerate}[label=(\alph*)]
\item Calculate Cassini's identity for these indices: $F_5 \times F_7 - F_6^2 = 5 \times 13 - 8^2$.
\item What happens to the slopes $\frac{F_{k-1}}{F_{k+1}}$ and $\frac{F_{k-2}}{F_k}$ as the Fibonacci indices grow larger? Does the gap become easier or harder for the human eye to detect?
\end{enumerate}
\end{enumerate}

\begin{walkbox}
Take an equation or rule of thumb you rely on frequently in your own work. Do you know the precise boundary conditions where it ceases to be true? What hidden assumption prevents it from failing?
\end{walkbox}

\newpage
